\documentclass[10pt,a4paper]{amsart}
\usepackage[margin=19mm]{geometry}
\usepackage{amsmath,amssymb,mathtools}
\usepackage[scr=rsfs,cal=euler]{mathalfa}
\usepackage{xfrac}
\usepackage[unicode,hidelinks,bookmarks=false]{hyperref}
\newcommand{\ie}{i.e.\ }
\newcommand{\cf}{cf.\ }
\newcommand{\resp}{respectively\ }
\newcommand{\Subsection}{\subsection}
\providecommand{\nobreakdash}{\nobreak\hbox{-}}
\setlength{\emergencystretch}{2em}
\multlinegap0pt
\usepackage{amssymb}
\usepackage{float}
\usepackage{mathtools}
\usepackage{xcolor}
\usepackage{algorithm}\usepackage{algpseudocode}
\usepackage{tikz}
\newtheorem{assumption}{Assumption}
\newtheorem{lemma}{Lemma}
\title{\huge A note on diffusive/random-walk behaviour in Metropolis--Hastings algorithms}
\author{Yuxin Liu, Peiyi Zhou, Samuel Livingstone}
\date{Source: arXiv:2603.07310v1 (CC BY 4.0)}
\begin{document}
\maketitle

\noindent\emph{Source and licence.} \huge A note on diffusive/random-walk behaviour in Metropolis--Hastings algorithms. Version arXiv:2603.07310v1 (CC BY 4.0), distributed by arXiv under the Creative Commons Attribution 4.0 International licence, http://creativecommons.org/licenses/by/4.0/, as recorded in the arXiv licence metadata for this exact version and shown on the abstract page https://arxiv.org/abs/2603.07310v1. Full licence notice: \texttt{licenses/arxiv-2603.07310-CC-BY-4.0.txt}.\par\smallskip

\noindent\textit{Editorial scope.} Counted unit: the article's lemma lem:upperbound\_polynomialtails (source0.tex lines 1260--1263). The article's complete original proof of it is delivered with it (source0.tex lines 1265--1304, one \texttt{proof} environment of 40 lines). No mathematics is written, repaired or re-derived: the statement and the proof are the article's own lines, in the article's own notation, and no retained line is substituted or re-flowed.  Retained uncounted as the source-local context the proof needs: lines 129--138; lines 516--522.  Nothing was omitted from any retained span, so the package carries no omission record.  No retained line contains a citation, so the wrapper carries no bibliography and no bibliography entry.  No retained line contains a figure environment, an image-inclusion command or drawing code, so no image is delivered and the package declares no asset dependency.  Editorial formatting only, in addition: an AMS \texttt{amsart} preamble that restates the article's own declarations and macros used by the retained lines, namely the environments \texttt{assumption}, \texttt{lemma} on separate counters, together with the standard packages those lines need.  The retained environments print in this document as Assumption 1, Lemma 1 in order of appearance, whereas the article prints the same items with the numbering of its own full document.  The complete native source of the article as distributed in the arXiv e-print for this exact version is bundled unchanged as \texttt{sources/195936/source0.tex}.
\begin{assumption}
    The distribution $\pi$ is defined on $\mathbb{R}^d$ and the following hold
    \begin{itemize}
        \item[(i)] $\pi$ has a density $\pi(x)$ that is upper semicontinuous and bounded away from 0 and $\infty$ on compact sets.
        \item[(ii)] $Q(x,\cdot)$ has a density $q(x,y)$, both $q(x,\cdot)$ and $q(\cdot,y)$ are lower semicontinuous for all $x$ and $y$, and there exists $\delta_q > 0$ such that for every $x \in \mathbb{R}^d$ \[
    \|x-y\| \leq \delta_q \implies q(x,y) > 0.
    \]
    \end{itemize}
    \label{ass:pi_q_regularities}
\end{assumption}

\begin{assumption}
\label{ass:heavy-tailed}
Assumption \ref{ass:pi_q_regularities}(i) holds and in addition there exist constants $r>0$, $K>0$, and $C_0>0$ such that for $|x|>K$,
$$
\pi(x) = \frac{C_0}{|x|^{1+r}}.
$$  
\end{assumption}

\begin{lemma}
\label{lem:upperbound_polynomialtails}
    If $\pi$ satisfies Assumption \ref{ass:heavy-tailed}, then the polynomial convergence rate of the guided walk Metropolis with finite variance proposal is at most $r$.
\end{lemma}

\begin{proof}
Suppose that the algorithm begins from $(x,p) = (0,1)$. For any $n>0$, we can choose $k>0$ such that $kn>K$. Define the set $A_n = \{x:x>kn\}$, and let $\mu$ be the joint invariant measure of $X,P$. For notational simplicity we write $P^n((x,p),A)$ for $P^n((x,p),(A,\{+1,-1\}))$, and note that $\mu((A,\{+1,-1\})) = \pi(A)$.  By the definition of total variation distance, we have
\begin{align}
    \|P^n((0,1),\cdot) - \mu\|_{TV} &\ge |P^n((0,1),A_n) - \pi(A_n)|.
\end{align}    
To bound $P^n((0,1),A_n)$, we observe that if the chain is never rejected in $n$ steps, it continuously moves in the positive direction. By a stochastic comparison argument the probability that the chain reaches $A_n$ is at most the probability of reaching $A_n$ without any rejections. Therefore,
\begin{equation*}
    P^n((0,1),A_n) \le \mathbb{P}\left(\sum_{i=1}^n |z_i|>kn\right), \quad \text{where } z_i \stackrel{\text{i.i.d.}}{\sim} N(0,\sigma^2) \text{for} \, i=1, \cdots ,n.
\end{equation*}
Since each $|z_i|$ follows a half-normal distribution with tail decay at least as fast as a Gaussian, $|z_i|$ is sub-Gaussian. Let $\mu = \mathbb{E}[|z_i|] = \sigma\sqrt{2/\pi}$ and $\sigma_H^2 = \text{Var}(|z_i|)$. Then for all $\lambda>0$,
\begin{equation*}
    \mathbb{E}\left[e^{\lambda(\sum_{i=1}^n |z_i| - n\mu)}\right]\le \exp\left(\frac{n\lambda^2\sigma_H^2}{2}\right).
\end{equation*}
Applying Markov's inequality to the exponential function, we obtain for all $\lambda>0$
\begin{equation}
\begin{aligned}
    \mathbb{P}\left(\sum_{i=1}^n|z_i| >kn\right) &= \mathbb{P}\left(\sum_{i=1}^n|z_i| - n\mu>(k-\mu)n\right)\\
    &= \mathbb{P}\left(\exp\left(\lambda \sum_{i=1}^n(|z_i| - \mu)\right) >\exp(\lambda(k-\mu)n)\right)\\
    &\le \frac{\mathbb{E}\left[\exp\left(\lambda \sum_{i=1}^n(|z_i| - \mu)\right)\right]}{\exp(\lambda(k-\mu)n)}\\
    &\le \exp\left(\frac{n\lambda^2\sigma_H^2}{2} - \lambda n (k-\mu)\right).
\end{aligned}
\end{equation}
Since this holds for all $\lambda>0$, we can minimize over $\lambda$ by choosing $\lambda^* = (k-\mu)/\sigma_H^2$, giving
\begin{align}
    P^n((0,1),A_n) \le  \mathbb{P}\left(\sum_{i=1}^n|z_i| >kn\right)\le \exp\left(-\frac{n(k-\mu)^2}{2\sigma_H^2}\right).
\end{align}
Since $kn>K$ by construction, for $x>kn$ we have $\pi(x) = C_0 x^{-(1+r)}$. Combining gives
\begin{equation}
    \pi(A_n) = \int_{kn}^\infty \frac{C_0}{x^{1+r}}  dx = \frac{C_0}{r(kn)^r} = \frac{C'}{n^r}, \quad \text{where } C' = C_0/(rk^r),
\end{equation}  
which implies that for all $n\ge n_0$, for suitably large $n_0$ and constant $k>0$
\begin{equation}
    |\pi(A_n) - P^n((0,1),A_n)| \ge \left|\frac{C'}{n^r} - \exp\left(-\frac{n(k-\mu)^2}{2\sigma_H^2}\right)\right| \ge \frac{c}{n^r}
\end{equation}
where $c>0$ is a constant. The last inequality holds because for any fixed $r>0$, we can choose large $k=k(r)$, and $C' = C_0/(rk^r):= C'(r)$ such that both relate to $r$ and $\exp(-n(k-\mu)^2/(2\sigma_H^2)) < C'/(2n^r)$ for all $n\ge n_0$ for suitably large $n_0$, which implies the bound with $c = C'/2$.  We have therefore established that
\begin{equation}
    \|P^n((0,1),\cdot) - \pi(\cdot)\|_{TV} \ge \frac{c}{n^r}, \quad n\ge n_0,
\end{equation}
from which the result follows.
\end{proof}

\end{document}
